% [原文 p.185]
% 承接 p.184 末句（C.4 节的证明）：
换言之，
\begin{equation}
  \eta(s)\psi\eta(s)^{-1}=s\psi\qquad(s\in S(\mathbb{H}),\ \psi\in\R^{4}) \tag{C.4.2}
\end{equation}
成立。$S(\mathbb{H})$ 的对合映射与 $\eta$ 并不相容，因为在偶数次部分 $C^{0}_{4}$ 上对合映射是恒等映射。

设 $\mathrm{Cliff}(\R^{4,0})$ 自然嵌入 $\mathrm{Cliff}(\R^{4,4})$，把所求的映射 $\nu$ 定义为
\[
  \nu(s,x,y,n)=(s,\ x+isy,\ \eta(s)n).
\]
由 $w$ 的定义可知
\[
  \eta(s)w=-\eta(-s),
\]
从而
\[
  \eta(-s)\tilde n=\eta(-s)\{-w\bar n\}=\eta(s)\bar n=\overline{\eta(s)n},
\]
故 $\nu$ 是实映射。由等式 (C.4.2) 可知
\[
  \eta(s)(x,y)n=(x+isy)\eta(s)n,
\]
故 $\nu$ 与模结构相容。这样便证明了图式 (C.4.1) 的存在性，引理 C.3.9 的证明完成。

这里给出的 $M(p,q)$ 与 $A(p,q)$ 从本文的观点看是自然的定义。但最好说明一下：用更具体或更经典的术语，它们对应着什么。为此注意，若 $M$ 是实 $C(\R^{p,q})$ 模，则可用
\[
  [x,y]m=xm+iym
\]
定义 $\R^{p+q}$ 在 $M$ 上的新作用 $[\ ,\ ]$。此时
\[
  [x,y]^{2}m=\{-\|x\|^{2}+\|y\|^{2}\}m.
\]
此外，关于对合，

% [原文 p.186]
有
\begin{align*}
  \overline{[x,y]m} &= x\bar m+i y\bar m \\
  &= x\bar m+i\,\overline{y}\bar m \qquad(\bar y=-y\text{ 故})\\
  &= [x,y]\bar m .
\end{align*}
因此 $M_{\R}$ 是二次曲面
\[
  Q(p,q)=\sum_{i=1}^{p}y_{i}^{2}-\sum_{j=1}^{q}x_{j}^{2}
\]
的 Clifford 代数 $C_{p,q}$ 在通常意义下的实模。容易看出这一手续可以倒过来走。因此 $M(p,q)$ 也可以定义为分次实 $C_{p,q}$ 模的 Grothendieck 群。由此，与 [4, §4, §5] 一样计算群 $A(p,q)$，并看出它们只依赖于 $p-q \pmod 8$，并不困难（亦见 [8]）。利用 [4, 推论 11.4] 的结果可知，
\[
  \alpha\colon A(p,q)\to KR^{p,q}(*)
\]
总是同构。细节留给读者。此时应当指出：这种双下标记法是受 Karoubi 的论文 [8] 启发而来的。

映射 $\alpha$ 还可以像 [4] 那样对主 $\mathrm{Spin}$ 丛一般地定义。于是，与 [4, 定理 12.3] 一样，可得 $\mathrm{Spin}$ 丛的 Thom 同构定理。把这些写成公式化陈述的事宜留给读者。

\section{与指标的关系}
\takeaway{本节把实结构推广到椭圆型算子：由算子 $P$ 的对合 $\overline{P}$ 导出表象的共轭关系 (C.5.1)，并说明实椭圆型算子的指数取值于 $KR$，在带参数的族情形下 $KR(Y)\to K(Y)$ 未必单射。}

记 $\hat\varphi$ 为函数 $\varphi$ 的 Fourier 变换，则
\[
  \widehat{\bar\varphi}(x)=\overline{\hat\varphi(-x)}
\]
成立。椭圆型微分算子 $P$ 的象征 $\sigma(P)$ 由 Fourier 变换定义 [9]，故

% [原文 p.187]
\[
  \sigma(\overline{P})(x,\xi)=\overline{\sigma(P)(x,-\xi)}
\]
可知。其中 $\overline{P}$ 是由
\[
  \overline{P}(\varphi)=\overline{P(\bar\varphi)}
\]
定义的算子。这里假定 $P$ 作用在函数上，以便 $\overline{P(\bar\varphi)}$ 有定义。更一般地，若 $X$ 是实可微流形，即带有可微对合 $x\mapsto\bar x$ 的可微流形，$E,F$ 是 $X$ 上的光滑实向量丛，则光滑截面空间 $\Gamma(E),\Gamma(F)$ 带有实结构，对任意线性算子
\[
  P\colon\Gamma(E)\to\Gamma(F),
\]
可按
\[
  \overline{P}(\varphi)=\overline{P(\bar\varphi)}
\]
定义 $\overline{P}\colon\Gamma(E)\to\Gamma(F)$。

若 $P$ 是椭圆型微分算子，则
\begin{equation}
  \sigma(\overline{P})(x,\xi)=\overline{\sigma(P)(\bar x,-\tau^{*}(\xi))} \tag{C.5.1}
\end{equation}
成立。当 $P=\overline{P}$ 时，把 $P$ 定义为实算子是自然的。若 $X$ 上的对合是恒等映射，这表示 $P$ 关于 $E,F$ 的实局部基具有实系数的微分算子。无论哪种情形，由 (C.5.1) 可知，实椭圆型算子的象征 $\sigma(P)$ 给出实向量丛的同构
\[
  \pi^{*}E\to\pi^{*}F,
\]
其中 $\pi\colon S(X)\to X$ 是余切球面丛的投影，$S(X)$ 上的对合由
\[
  (x,\xi)\to(\bar x,-\tau^{*}(\xi))
\]
定义。当 $X$ 上的对合 $\tau$ 是恒等映射时，$S(X)$ 上的对合并不是恒等映射，而是在每条纤维上成为对合映射，请注意这一点。

% [原文 p.188]
这正是研究实椭圆型算子需要 $KR$ 理论的基本理由。事实上，三元组
\[
  (\pi^{*}E,\ \pi^{*}F,\ \sigma(P))
\]
按通常的方式定义一个元素
\[
  [\sigma(P)]\in KR\bigl(B(X),S(X)\bigr).
\]
其中 $T^{*}X$ 的单位球丛 $B(X)$ 也具有与 $S(X)$ 相同的实结构\footnote{当然，这一构造对一切积分算子（或拟微分算子）都可推广。}。

实椭圆型算子的核与余核带有自然的实结构，因此指数自然是 $KR(*)$ 的元素。当然，$KR(*)\to K(*)$ 是同构，所以这样精化后的实指数并没有直接的好处。但若考虑实椭圆型算子的族，情况就不同了。设 $Y$ 为参数空间，则实指数是 $KR(Y)$ 的元素，而 $KR(Y)\to K(Y)$ 一般未必是单射。

这些事情都可自然推广到实椭圆型复形 [9]。特别有趣的是实代数簇上的 Dolbeault 复形。它是实椭圆型复形，因为正则映射 $\tau\colon X\to\bar X$ 把 $\bar X$ 的 Dolbeault 复形映入 $X$ 的 Dolbeault 复形。若 $X$ 的层上同调满足 $q\ge 1$ 时 $H^{q}(X,\mathcal{O})=0$，且 $H^{0}(X,\mathcal{O})\cong\C$，则 Dolbeault 复形的指数即 Euler 示性是 $1$。利用这一事实可以证明下述命题。

\begin{proposition*}{C.5.2}
  设 $f\colon X\to Y$ 是实代数簇的纤维丛，其纤维 $F$ 满足
  \[
    H^{q}(F,\mathcal{O})=0\qquad(q\ge 1),\qquad H^{0}(F,\mathcal{O})\cong\C,
  \]
  则存在同态
  \[
    f_{*}\colon KR(X)\to KR(Y),
  \]
  它是 $f^{*}\colon KR(Y)\to KR(X)$ 的左逆。
\end{proposition*}

% [原文 p.189]
这里无法给出它的证明，但看一个特殊情形。取 $F$ 为实线性代数群的（紧）齐性空间，$X=Y\times F$。例如，$F$ 可以取作复二次曲面，就像证明 (C.3.13) 时那样；也可以取 $F$ 为 $\mathrm{SO}(2n)/\mathrm{U}(n)$，或 $\mathrm{SO}(2n)$ 的旗流形 $\mathrm{SO}(2n)/T^{n}$。这些空间可用于证明正交丛的分裂原理。此时，看出实空间
\[
  \bigl(\mathrm{SO}(2n)/\mathrm{U}(n)\bigr)\times\R^{0,2n}
\]
具有实向量丛的结构，是很重要的。$\mathrm{SO}(2n)/\mathrm{U}(n)$ 的点给出 $\R^{2n}$ 的复结构，共轭点给出共轭复结构。取 $n=2$ 时，这本质上就是 C.3 节中由定理 C.2.1 导出正交周期性时所做之事\footnote{命题 C.3.1 中用了 $3$ 维球面 $S^{4,0}$，同样也可以用 $2$ 维球面 $S^{3,0}$，这对应于使用 $\mathrm{SO}(4)/\mathrm{U}(2)$。}。

\section*{参考文献}
\addcontentsline{toc}{section}{参考文献}
\begingroup
\small
\noindent [1] D. W. Anderson, `A new cohomology theory', Thesis, Univ. of California, Berkeley (1964).

\noindent [2] M. F. Atiyah, `Lectures on K-theory' (mimeographed notes, Harvard 1965).

% [原文 p.190]
\noindent [3] M. F. Atiyah and R. Bott, `On the periodicity theorem for complex vector bundles', \textit{Acta Math.} \textbf{112} (1964) 229--247.

\noindent [4] M. F. Atiyah, R. Bott, and A. Shapiro, `Clifford modules', \textit{Topology} \textbf{3}, \textit{Suppl. 1} (1964) 3--38.

\noindent [5] M. F. Atiyah and F. Hirzebruch, `Vector bundles and homogeneous spaces', Proc. of Symp. in Pure Math. 3 American Math. Soc. (1961) 7--38.

\noindent [6] M. F. Atiyah and G. B. Segal, `Equivariant K-theory' (Lecture notes, Oxford 1965).

\noindent [7] P. S. Green, `A cohomology theory based upon self-conjugacies of complex vector bundles', \textit{Bull. Amer. Math. Soc.} \textbf{70}(4) (1964) 522--524.

\noindent [8] M. Karoubi, `Algèbres de Clifford et K-théorie', \textit{Annales scientifiques de l'École Normale Supérieure}, Série 4, \textbf{1}(2) (1968) 161--270.

\noindent [9] R. Palais, \textit{Seminar on the Atiyah--Singer index theorem}, Ann. of Math. Studies 57 (Princeton University Press, Princeton 1965).
\endgroup

\chapterreview{%
附录 C 是全书最晚成文、也最“技术外挂”的一篇：它把前面用 $K$ 理论处理复向量丛的全套机器，搬到实向量丛与带对合的空间上，从而把实 $K$ 理论（$KR$ 理论）安放进同一个 Grothendieck 框架。读它的动机应当放在 Atiyah--Singer 指标定理的实形式里：实椭圆型算子自然给出 $KR$ 中的元素，而不再是 $K$ 中的元素。

知识脉络上，本篇从实空间的范畴出发（C.1），把“实结构”表述为空间与丛上的对合；C.2 重证周期性，说明 Bott 周期在实的情形呈现 $8$ 周期；C.3 用系数理论把 $KR^{p,q}$ 组织成双下标理论，并借 Clifford 代数给出代数模型；C.4 把这套双下标记法与 $C_{p,q}$ 模的 Grothendieck 群显式对应；C.5 才回到算子，把实结构与椭圆型算子的象征、指数接上。

核心结论有三条：其一，实 $KR$ 理论满足 $8$ 周期，且 $KR^{p,q}$ 只依赖 $p-q \pmod 8$；其二，$M(p,q)$ 与 $A(p,q)$ 可以用 Clifford 模的 Grothendieck 群具体刻画，于是 $\alpha\colon A(p,q)\to KR^{p,q}(*)$ 恒为同构；其三，对实椭圆型算子，象征给出 $S(X)$ 上实向量丛的同构，指数落在 $KR(*)$ 里，而对带参数的族，$KR(Y)\to K(Y)$ 未必单射，这正是实指标比复指标更精细的地方。

辨析与提醒：第一，本篇的对合有两层——空间 $X$ 上的 $\tau$ 与余切球面丛 $S(X)$ 上的 $(x,\xi)\mapsto(\bar x,-\tau^{*}(\xi))$，二者不要混为一谈；当 $\tau=\mathrm{id}$ 时后者在每条纤维上是 $-1$，这是“实算子”与“复算子”分野的根源。第二，$KR^{p,q}$ 的上下标是双下标，与本篇之前 $K^{n}(X)$ 的单下标不同，引用时务必写全。第三，C.5 中 $P=\overline{P}$ 才称实算子，作者把这一条件写成自然定义，但读者应自行验证它与“实局部基下实系数”的等价性。第四，本篇大量把细节留给读者（例如 $\alpha$ 的同型性、Spin 丛的 Thom 同构、命题 C.5.2 的证明），这是原文风格，本译本未作补写。

评注：若只读正文三章，$K$ 理论看起来是复向量丛的上同调论；附录 C 提醒我们，它其实是一个可以对任意“带对合”的范畴作系数的理论。把 C.3 的系数理论与 C.5 的指标放在一起看，才能理解为什么实指标定理必须用 $KR$ 而不是 $K$ 来陈述。}
% CHUNK appC3 DONE
